\section{Experiments}
\label{sec::experiment}
\begin{table*}[t]
\centering

% --- 定义新样式 ---
\newcommand{\res}[2]{%
  #1{\scriptsize\,\ensuremath{\pm}\,#2}%
}
\newcommand{\best}[2]{%
  \textbf{#1{\scriptsize\,\ensuremath{\pm}\,#2}}%
}
\newcommand{\second}[2]{%
  \underline{#1{\scriptsize\,\ensuremath{\pm}\,#2}}%
}

\newcommand{\head}[1]{\mbox{\textbf{#1}}}
\newcolumntype{Y}{>{\centering\arraybackslash}X}

\renewcommand{\arraystretch}{0.95} 
\setlength{\tabcolsep}{1pt}

\resizebox{\linewidth}{!}{
\begin{tabularx}{1.5\linewidth}{c c | YYYYYYYYYYY} 
\toprule

\textbf{Datasets} & \textbf{Metric}
& \head{pcaReduce} & \head{DEC} & \head{contrastive-sc} & \head{scNAME} & \head{scDeepCluster}
& \head{scDSC} & \head{AttentionAE-sc} & \head{scGNN} & \head{scCDCG} & \head{scSiameseClu} & \head{Ours} \\
\midrule

\multirow{3}{*}{\shortstack{\textbf{Muraro Human}\\\textbf{Pancreas cells}}}
& ACC & \res{50.56}{3.4} & \res{72.22}{5.5} & \res{81.85}{4.6} & \res{80.36}{7.9} & \res{74.70}{2.7} & \res{79.42}{1.4} & \res{93.84}{2.7} & \res{79.32}{4.0} & \res{92.65}{1.9} & \second{94.95}{0.1} & \best{96.02}{0.5} \\
& NMI & \res{54.83}{1.7} & \res{76.29}{2.7} & \res{77.73}{3.3} & \res{79.12}{2.7} & \res{79.32}{0.3} & \res{75.89}{0.6} & \second{87.80}{2.3} & \res{78.76}{5.6} & \res{86.81}{1.0} & \res{86.19}{0.3} & \best{89.15}{1.1} \\
& ARI & \res{36.61}{2.3} & \res{61.76}{5.6} & \res{75.22}{10.1} & \res{76.64}{10.4} & \res{64.59}{2.5} & \res{75.42}{1.2} & \res{90.74}{4.1} & \res{78.87}{3.2} & \res{91.37}{1.2} & \second{91.59}{0.3} & \best{93.10}{0.9} \\
\midrule

\multirow{3}{*}{\shortstack{\textbf{Human Pancreas}\\\textbf{cells 1}}}
& ACC & \res{23.63}{0.2} & \res{40.64}{1.9} & \res{69.67}{7.8} & \res{84.20}{6.0} & \res{66.39}{3.8} & \res{73.10}{2.3} & \res{81.56}{1.8} & \res{55.10}{2.1} & \res{92.15}{0.8} & \second{95.12}{0.4} & \best{97.21}{0.3} \\
& NMI & \res{42.58}{0.4} & \res{62.35}{0.6} & \res{69.21}{3.3} & \res{65.74}{4.0} & \res{69.42}{12.5} & \res{60.50}{4.2} & \res{82.62}{3.0} & \res{64.03}{0.9} & \res{86.35}{0.9} & \second{89.50}{0.6} & \best{93.23}{0.2} \\
& ARI & \res{1.98}{0.5} & \res{26.07}{1.3} & \res{52.00}{8.7} & \res{63.31}{7.4} & \res{47.92}{3.0} & \res{69.81}{1.6} & \res{71.84}{8.5} & \res{38.32}{0.8} & \res{92.83}{0.6} & \second{93.40}{0.5} & \best{96.61}{1.5} \\
\midrule

\multirow{3}{*}{\shortstack{\textbf{Human Pancreas}\\\textbf{cells 2}}}
& ACC & \res{37.56}{0.4} & \res{45.26}{1.8} & \res{60.01}{5.7} & \res{63.33}{2.0} & \res{62.88}{5.5} & \res{80.40}{4.0} & \res{80.40}{4.0} & \res{57.62}{2.3} & \res{84.66}{1.1} & \second{89.45}{0.9} & \best{91.13}{0.9} \\
& NMI & \res{50.81}{0.3} & \res{66.32}{2.1} & \res{62.25}{2.3} & \res{58.05}{1.6} & \res{73.97}{2.2} & \res{79.44}{1.7} & \res{82.16}{10.9} & \res{79.32}{2.3} & \res{83.44}{2.6} & \second{86.80}{1.2} & \best{90.63}{1.2} \\
& ARI & \res{19.23}{0.5} & \res{39.14}{1.8} & \res{47.10}{5.3} & \res{32.12}{1.5} & \res{64.56}{7.5} & \res{82.85}{6.5} & \res{73.20}{10.5} & \res{79.96}{1.7} & \res{85.30}{1.6} & \second{89.20}{1.1} & \best{92.38}{1.4} \\
\midrule

\multirow{3}{*}{\shortstack{\textbf{Human Pancreas}\\\textbf{cells 3}}}
& ACC & \res{35.92}{0.4} & \res{43.56}{1.5} & \res{57.47}{5.0} & \res{83.20}{8.7} & \res{73.65}{3.1} & \res{80.12}{10.0} & \second{90.96}{2.4} & \res{67.94}{4.6} & \res{88.04}{0.3} & \res{89.45}{0.5} & \best{95.31}{1.4} \\
& NMI & \res{51.22}{0.3} & \res{64.51}{0.8} & \res{66.84}{2.3} & \res{80.56}{5.6} & \res{75.02}{1.2} & \res{75.87}{8.4} & \second{87.00}{1.8} & \res{62.60}{1.8} & \res{82.21}{1.1} & \res{86.20}{0.6} & \best{91.37}{0.3} \\
& ARI & \res{26.40}{4.5} & \res{42.73}{1.0} & \res{47.90}{3.4} & \res{81.08}{5.7} & \res{63.31}{2.3} & \res{82.85}{8.5} & \res{89.78}{2.5} & \res{58.41}{1.9} & \res{90.78}{0.6} & \second{92.34}{0.8} & \best{95.25}{0.5} \\
\midrule

\multirow{3}{*}{\shortstack{\textbf{Mouse Pancreas}\\\textbf{cells 2}}}
& ACC & \res{28.97}{0.3} & \res{34.19}{2.3} & \res{59.55}{5.7} & \res{82.36}{8.0} & \res{69.91}{2.9} & \res{77.35}{2.3} & \res{81.92}{6.5} & \res{43.03}{3.0} & \second{93.97}{0.2} & \res{92.45}{0.4} & \best{95.11}{1.2} \\
& NMI & \res{48.93}{0.1} & \res{54.08}{0.9} & \res{62.34}{2.4} & \res{78.00}{3.9} & \res{59.58}{2.6} & \res{62.50}{0.4} & \res{81.88}{5.1} & \res{55.87}{1.9} & \second{88.02}{0.2} & \res{86.50}{0.5} & \best{88.59}{0.6} \\
& ARI & \res{12.80}{0.1} & \res{18.54}{1.1} & \res{49.75}{8.5} & \res{76.56}{11.4} & \res{60.14}{3.6} & \res{65.81}{1.4} & \res{82.14}{7.8} & \res{26.40}{2.5} & \second{92.61}{0.3} & \res{90.80}{0.6} & \best{93.43}{0.8} \\
\midrule

\multirow{3}{*}{\textbf{CITE-CMBC}}
& ACC & \res{28.19}{0.1} & \res{41.88}{5.4} & \res{52.70}{5.0} & \res{68.06}{13.3} & \res{70.80}{2.5} & \res{68.79}{0.9} & \res{61.18}{7.9} & \res{66.71}{4.0} & \res{71.45}{1.8} & \second{74.70}{0.5} & \best{79.20}{0.5} \\
& NMI & \res{28.36}{0.2} & \res{46.87}{9.2} & \res{64.83}{0.8} & \res{63.23}{13.2} & \res{72.53}{0.7} & \res{64.21}{3.3} & \res{62.02}{9.6} & \res{61.70}{3.1} & \second{74.77}{1.7} & \res{68.71}{0.7} & \best{76.50}{0.4} \\
& ARI & \res{5.10}{0.1} & \res{23.55}{10.3} & \res{47.88}{3.7} & \res{58.99}{21.3} & \res{56.23}{4.1} & \res{52.51}{2.2} & \res{44.20}{11.6} & \res{61.20}{3.1} & \res{61.46}{1.4} & \second{67.56}{1.1} & \best{69.80}{0.5} \\
\midrule

\multirow{3}{*}{\shortstack{\textbf{Human Liver}\\\textbf{cells}}}
& ACC & \res{34.92}{0.2} & \res{46.32}{5.5} & \res{79.31}{2.1} & \res{74.55}{5.3} & \res{63.44}{0.7} & \res{70.90}{2.2} & \res{79.46}{6.7} & \res{68.65}{3.2} & \res{75.34}{1.7} & \second{88.33}{1.7} & \best{94.50}{0.9} \\
& NMI & \res{32.07}{0.6} & \res{52.24}{5.0} & \res{85.11}{0.1} & \res{77.42}{12.0} & \res{74.60}{2.5} & \res{71.63}{2.8} & \res{82.86}{4.5} & \res{62.33}{2.1} & \res{79.34}{2.6} & \second{88.82}{1.2} & \best{90.15}{1.1} \\
& ARI & \res{6.34}{14.4} & \res{31.04}{4.7} & \res{89.96}{0.1} & \res{79.91}{2.9} & \res{58.58}{10.1} & \res{75.34}{3.6} & \res{79.16}{11.8} & \res{65.41}{1.3} & \res{81.26}{2.7} & \second{91.90}{0.5} & \best{94.10}{0.5} \\
\bottomrule
\end{tabularx}
}
\caption{Clustering performance across seven benchmark datasets. The results are reported as mean $\pm$ standard deviation, where the best results are highlighted in \textbf{bold} and the second-best are \underline{underlined}.}
\label{tab:main_result}
\end{table*}
\begin{figure*}[t]
    \centering
    \begin{subfigure}[b]{0.14\linewidth}
        \includegraphics[width=\linewidth]{figs/UMAP/Sonya_HumanLiver_counts_top5000_scDeepcluter_umap_basic.pdf}
        \caption{scDeepCluster}
    \end{subfigure}
    \hfill
    \begin{subfigure}[b]{0.16\linewidth}
        \includegraphics[width=\linewidth]{figs/UMAP/Sonya_HumanLiver_counts_top5000_scDSC_umap_basic.pdf}
        \caption{scDSC}
    \end{subfigure}
    % \hfill
    % \begin{subfigure}[b]{0.16\linewidth}
    %     \includegraphics[width=\linewidth]{figs/UMAP/Sonya_HumanLiver_counts_top5000_ours_umap_AttentionAE.pdf}
    %     \caption{AttentionAE-sc}
    % \end{subfigure}
    \hfill
    \begin{subfigure}[b]{0.16\linewidth}
        \includegraphics[width=\linewidth]{figs/UMAP/Sonya_HumanLiver_counts_top5000_scGNN_umap_basic.pdf}
        \caption{scGNN}
    \end{subfigure}
    \hfill
    \begin{subfigure}[b]{0.16\linewidth}
        \includegraphics[width=\linewidth]{figs/UMAP/Sonya_HumanLiver_counts_top5000_scCDCG_umap_basic.pdf}
        \caption{scCDCG}
    \end{subfigure}
    \hfill
    \begin{subfigure}[b]{0.16\linewidth}
        \includegraphics[width=\linewidth]{figs/UMAP/Sonya_HumanLiver_counts_top5000_ours_umap_scSiameseClu.png}
        \caption{scSiameseClu}
    \end{subfigure}
    \hfill
    \begin{subfigure}[b]{0.16\linewidth}
        \includegraphics[width=\linewidth]{figs/UMAP/Sonya_HumanLiver_counts_top5000_ours_umap_basic.pdf}
        \caption{\method{}}
    \end{subfigure}
    \caption{UMAP visualizations of \method{} and five baseline methods on the Human Liver cells dataset. Each point represents a cell, and colors indicate the predicted cell type.}

    \label{fig:UMAP_comparison}
\end{figure*}

\subsection{Experimental Setup}


\smallskip\noindent\textbf{Benchmark Datasets}. To evaluate the robustness of our model for clustering across species and diverse biological contexts, we conduct experiments on seven public single-cell benchmark datasets, covering both human and mouse samples and including pancreas, liver, and peripheral immune cell populations. 
In terms of data scale, the number of cells ranges from 1{,}064 to 8{,}617, the number of genes ranges from 2{,}000 to 20{,}125, and the expression matrices exhibit sparsity rates of approximately 73.02\%--93.26\%. To standardize the input dimensionality and focus on more informative features, we select either 1{,}500 or 2{,}000 highly variable genes (HVGs) for each dataset as the model input. Detailed statistics are summarized in the Table~\ref{tab:dataset_summary} of Appendix~\ref{sec:setup}.

\smallskip\noindent\textbf{Baselines}. To highlight the advantages of \method{}, we compare it with ten representative clustering approaches. Specifically, pcaReduce~\cite{vzurauskiene2016pcareduce} is included as a representative traditional clustering method. DEC~\cite{xie-et-al:dec}, contrastive-sc ~\cite{ciortan2021contrastive}, scNAME~\cite{wan-et-al:scname}, scDeepCluster~\cite{tian-et-al:model-based-scrnaseq} and AttentionAE-sc~\cite{li-et-al:attention-scrnaseq} are deep learning--based clustering methods, which mainly learn representations via autoencoders or contrastive learning and then perform clustering. In addition, scGNN~\cite{wang2021scgnn}, scDSC~\cite{gan2022deep},  scCDCG~\cite{xu2024sccdcg} and scSiamese~\cite{xu2025scsiameseclu} are deep graph clustering methods that incorporate GNNs and structural information for representation learning and clustering. Detailed experimental settings and implementation details are provided in the Appendix~\ref{sec:setup}.

\smallskip\noindent\textbf{Evaluation Metrics}. To evaluate the effectiveness of the proposed method, we assess clustering performance using three widely adopted metrics: clustering accuracy (ACC), normalized mutual information (NMI)~\cite{strehl-ghosh:cluster-ensembles}, and adjusted Rand index (ARI)~\cite{vinh-et-al:ari}.

\begin{figure*}[t]
    \centering
    % \begin{subfigure}[b]{0.19\linewidth}
    %     \includegraphics[width=\linewidth]{figs/DEG_overlap_heatmap/cell_type_annotation_heatmap_scDeepCluster.pdf}
    %     \caption{scDeepCluster}
    % \end{subfigure}\hfill
    \begin{subfigure}[b]{0.19\linewidth}
        \includegraphics[width=\linewidth]{figs/DEG_overlap_heatmap/cell_type_annotation_heatmap_scDSC.pdf}
        \caption{scDSC}
    \end{subfigure}\hfill
    \begin{subfigure}[b]{0.19\linewidth}
        \includegraphics[width=\linewidth]{figs/DEG_overlap_heatmap/cell_type_annotation_heatmap_scGNN.pdf}
        \caption{scGNN}
    \end{subfigure}\hfill
    \begin{subfigure}[b]{0.19\linewidth}
        \includegraphics[width=\linewidth]{figs/DEG_overlap_heatmap/cell_type_annotation_heatmap_scCDCG.pdf}
        \caption{scCDCG}
    \end{subfigure}\hfill
    \begin{subfigure}[b]{0.19\linewidth}
        \includegraphics[width=\linewidth]{figs/DEG_overlap_heatmap/cell_type_annotation_heatmap_scSiameseclu.pdf}
        \caption{scSiameseClu}
    \end{subfigure}\hfill
    \begin{subfigure}[b]{0.19\linewidth}
        \includegraphics[width=\linewidth]{figs/DEG_overlap_heatmap/cell_type_annotation_heatmap_ours.pdf}
        \caption{\method{}}
    \end{subfigure}
    \caption{Cell type annotation. Heatmap of overlap between the top 100 DEGs in clusters detected by five methods and gold standard cell types (Cell types are abbreviated as PP, A, Ductal, B, Acinar, Epsilon, Stellate, D, and Endo).}
    \label{fig:cell_type_annotation_heatmap}
\end{figure*}

\begin{figure*}[t]
  \centering
  % --- 左图：Gold Standard Ranking ---
  \begin{subfigure}[t]{0.48\textwidth}
    \centering
    \includegraphics[width=\linewidth]{figs/Marker_Gene_Identification/deg_ranking_gold_standard.pdf}
    \caption{Gold Standard}
    \label{fig:deg_gold_10}
  \end{subfigure}
  \hfill
  % --- 右图：Model Prediction Ranking ---
  \begin{subfigure}[t]{0.48\textwidth}
    \centering
    \includegraphics[width=\linewidth]{figs/Marker_Gene_Identification/deg_ranking_model_prediction.pdf}
    \caption{\method{}}
    \label{fig:deg_ours_10}
  \end{subfigure}

  \caption{Top-5 DEGs Ranking. Comparison of the top 5 DEGs identified by the gold standard (a) and \method{} (b) for different cell types.}
  \label{fig:deg_ranking_comparison}
\end{figure*}

\subsection{Overall Performance}

\smallskip\noindent\textbf{Quantitative Analysis}. Table~\ref{tab:main_result} compares the clustering performance of \method{} with ten competitive baselines on seven benchmark datasets. From the results, we make the following observations. 
(1) \method{} achieves the best performance across all three metrics on all datasets, demonstrating its strong capability in handling high-dimensional and highly sparse scRNA-seq data. On average, compared with the second-best method, \method{} yields significant improvements of \textbf{3.00\%}, \textbf{2.42\%}, and \textbf{2.30\%} in terms of ACC, NMI, and ARI, respectively, validating the effectiveness and generalizability of our approach.
(2) Compared with deep clustering methods that rely solely on gene expression features, our \method{} consistently delivers substantial gains, highlighting the benefit of jointly leveraging expression features and cell--cell structural information.
(3) Although GNN-based methods like scGNN incorporate cell graphs, they predominantly rely on local neighborhood aggregation, which renders them vulnerable to over-smoothing and limits their ability to capture global dependencies. \method{} overcomes this by encoding complex structural information via the Siamese Graph Transformer, yielding more discriminative representations. Analysis of runtime is provided in Appendix~\ref{sec:runtime}.

\smallskip\noindent\textbf{Visualization}. To further investigate the model's ability to capture the underlying structure of cells in the embedding space, we apply UMAP~\cite{mcinnes2020umap} to project the learned representations of the Human Liver dataset into two dimensions, as shown in Figure~\ref{fig:UMAP_comparison}. We observe that our \method{} produces highly compact clusters with clear inter-cluster boundaries, faithfully reflecting the complex underlying tissue structure and effectively preserving the intrinsic heterogeneity among different cell types. In contrast, methods such as scDSC often yield projections with noticeable subtype mixing or diffuse distributions, making it difficult to separate rare cell populations from major groups.

\subsection{Biological Analysis}
% \smallskip\noindent\textbf{Cell Type Annotation}. Differentially expressed genes (DEGs) provide critical evidence for cell-type annotation. Using Seurat's "FindAllMarkers"~\cite{butler2018integrating}, we extracted the top-100 DEGs for each cluster on the Muraro Human Pancreas dataset and computed the overlap ratio with the DEG lists of the corresponding ground-truth cell types. As shown in Figure~\ref{fig:cell_type_annotation_heatmap}, The high-overlap regions of \method{} are mainly concentrated along the diagonal, surpassing 90\% overlap in most clusters, indicating that our clusters can be stably mapped to cell types with minimal confusion. In contrast, other methods exhibit lower overlap ratios, suggesting stronger inter-cluster confusion and less stable annotations.

% \smallskip\noindent\textbf{Marker Gene Identification}.
% To further investigate the biological fidelity of the learned representations, we focused on the specific marker genes associated with each identified population. First, as illustrated in Figure~\ref{fig:deg_ranking_comparison}, we compared the rankings of the top 10 DEGs between our \method{} and the Gold Standard. The remarkable consistency in the top-ranked genes confirms that our model captures the correct dominant gene signatures for distinct cell clusters. Second, to inspect the expression specificity, we visualized the top 3 marker genes using dot plots in Figure~\ref{fig:dotplot_comparison}. As observed, the expression patterns inferred by \method{} are nearly indistinguishable from the ground truth.  These results demonstrate that \method{} can effectively uncover both rare and major biologically meaningful subpopulations and their signature marker genes without supervision.
\smallskip\noindent\textbf{Differential Gene Expression Analysis}. 
Differentially expressed genes (DEGs) serve as a critical metric to validate the biological identity of learned clusters. Using Seurat's ``FindAllMarkers''~\cite{butler2018integrating}, we extract the top-100 DEGs for each cluster on the Muraro Human Pancreas dataset and quantitatively compute the overlap ratio with the corresponding ground-truth cell types. As shown in Figure~\ref{fig:cell_type_annotation_heatmap}, the high-overlap regions of \method{} are predominantly concentrated along the diagonal, consistently surpassing 90\% overlap in most clusters. This indicates that our clusters can be stably mapped to specific cell types with minimal confusion. In stark contrast, baseline methods frequently exhibit lower overlap ratios due to ambiguity and inter-cluster confusion. To further investigate the biological fidelity, we compare the rankings of the top 5 DEGs between \method{} and the Gold Standard as illustrated in Figure~\ref{fig:deg_ranking_comparison}. Beyond mere overlap, we observe remarkable consistency in the dominant markers. For instance, in the cluster corresponding to Pancreatic PP cells (Cluster 7), our \method{} identifies \textit{PPY} and \textit{ETV1} as the top-two markers, perfectly matching the Gold Standard. This precise alignment in both gene identity and ranking confidence demonstrates that our \method{} does not merely capture global similarities but accurately prioritizes the key marker genes essential for distinct cell identification.
% \input{figs/Marker_Gene_Identification}

\smallskip\noindent\textbf{KEGG Pathway Enrichment Analysis}. To validate the biological interpretability of the clustering results, we perform KEGG pathway enrichment analysis on representative clusters. As shown in Figure~\ref{fig:kegg_bar}, our \method{} successfully resolves distinct functional subpopulations in pancreatic tissue. Cluster 3 is significantly enriched in ``Insulin secretion'' and critical diabetes-related pathways including MODY, while also exhibiting neuroendocrine signatures such as the ``Synaptic vesicle cycle'', confirming its biological identity as pancreatic $\beta$ cells. Cluster 4 shows strong enrichment in ``Ribosome'' and protein synthesis-related pathways, highlighting the high protein synthesis activity characteristic of acinar cells. Finally, Cluster 5 is enriched in pathways such as ``Phagosome'' and ``Lysosome'', indicating the presence of tissue-resident immune cells such as macrophages. These results demonstrate that our \method{} can accurately identify KEGG pathways and distinct functional phenotypes.

Furthermore, we provide trajectory inference analysis in Appendix~\ref{sec:trajectory} to validate preserved developmental manifolds.

\begin{figure*}[t]
    \centering
    \begin{subfigure}[b]{0.32\linewidth}
        \includegraphics[width=\linewidth]{figs/KEGG/Cluster_3_KEGG_enrichment.pdf}
        \caption{Cluster 3}
    \end{subfigure}\hfill
    \begin{subfigure}[b]{0.32\linewidth}
        \includegraphics[width=\linewidth]{figs/KEGG/Cluster_4_KEGG_enrichment.pdf}
        \caption{Cluster 4}
    \end{subfigure}\hfill
    \begin{subfigure}[b]{0.32\linewidth}
        \includegraphics[width=\linewidth]{figs/KEGG/Cluster_5_KEGG_enrichment.pdf}
        \caption{Cluster 5}
    \end{subfigure}
    \caption{Top 5 KEGG pathways for each cluster (bar graph shows $-\log_{10}(q)$).}
    \label{fig:kegg_bar}
\end{figure*}
 % \begin{figure*}[t]
 %      \centering
 %      \begin{subfigure}[b]{0.28\linewidth}
 %          \includegraphics[width=\linewidth]{IJCAI_stGTN/figs/sensitivity/sim_3dbar_ACC_clean.pdf}
 %          \caption{ACC}
 %      \end{subfigure}
 %      \begin{subfigure}[b]{0.28\linewidth}
 %          \includegraphics[width=\linewidth]{IJCAI_stGTN/figs/sensitivity/sim_3dbar_ARI_clean.pdf}
 %          \caption{NMI}
 %      \end{subfigure}
 %      \begin{subfigure}[b]{0.28\linewidth}
 %          \includegraphics[width=\linewidth]{IJCAI_stGTN/figs/sensitivity/sim_3dbar_NMI_clean.pdf}
 %          \caption{ARI}
 %      \end{subfigure}
 %      \caption{ The impact of two key hyperparameters on clustering performance.}
 %      \label{fig:sensitivity}
 %  \end{figure*}

 \begin{figure}[t]
    \centering
    \begin{subfigure}[b]{0.48\linewidth}
        \includegraphics[width=\linewidth]{figs/sensitivity/sim_3dbar_ACC_clean.pdf}
        \caption{ACC}
    \end{subfigure}
    \hfill
    \begin{subfigure}[b]{0.48\linewidth}
        \includegraphics[width=\linewidth]{figs/sensitivity/sim_3dbar_NMI_clean.pdf}
        \caption{NMI}
    \end{subfigure}
    \caption{The impact of two key hyperparameters ($K$ and $L$) on clustering performance in terms of ACC and NMI.}
    \label{fig:sensitivity}
\end{figure}
\subsection{Ablation Study}

\smallskip\noindent\textbf{Effectiveness of Major Components.}
We evaluate the contribution of each loss term on the Muraro dataset. Table~\ref{tab:main_ablation} indicates that the exclusion of any component leads to a sharp performance degradation of over 10\% in ACC, highlighting the non-redundancy of each module. Specifically, the absence of the optimal transport clustering loss $\mathcal{L}_{\text{clu}}$ results in the most severe drop, underscoring its critical role in guiding partition assignments. Similarly, removing the correlation term $\mathcal{L}_{\text{cor}}$ or the reconstruction objectives $\mathcal{L}_{\text{ZINB}}$ and $\mathcal{L}_{\text{rec}}$ significantly impairs the capability of the model to align views and preserve topology. Ultimately, the full \method{} achieves superior performance, validating the synergy of jointly optimizing for structural, feature, and distribution constraints. Dual-component ablation analysis is provided in Appendix~\ref{sec:extended_ablation}.
\begin{table}[t]
\centering


\small
\begin{tabular}{cccc|ccc}
\toprule
\boldmath$\mathcal{L}_{\text{clu}}$ & \boldmath$\mathcal{L}_{\text{cor}}$ & \boldmath$\mathcal{L}_{\text{ZINB}}$ & \boldmath$\mathcal{L}_{\text{rec}}$ & \textbf{ACC} & \textbf{NMI} & \textbf{ARI} \\
\midrule
$\checkmark$ & $\checkmark$ & $\checkmark$ & $\times$      & 84.57 & 78.93 & 76.08 \\
$\checkmark$ & $\checkmark$ & $\times$      & $\checkmark$ & 84.18 & 78.36 & 75.62 \\
$\checkmark$ & $\times$      & $\checkmark$ & $\checkmark$ & 83.47 & 77.58 & 74.36 \\
$\times$      & $\checkmark$ & $\checkmark$ & $\checkmark$ & 81.24 & 74.79 & 70.68 \\
\midrule
$\checkmark$ & $\checkmark$ & $\checkmark$ & $\checkmark$ & \textbf{96.02} & \textbf{89.15} & \textbf{93.10} \\
\bottomrule
\end{tabular}
\caption{Ablation experiment results of major loss terms}
\label{tab:main_ablation}
\end{table}

\smallskip\noindent\textbf{Impact of Embedding Components.} We further dissect the embedding module to assess the contributions of Gene Expression (GE), Position (Pos), and Shortest-Path (SP) embeddings in Table~\ref{tab:ablation_embed}. Results indicate that while GE provides the foundational semantics, integrating it with structural priors yields massive gains. Notably, adding Pos alone boosts ACC by over 12\% compared to the GE baseline, and the full combination achieves peak performance, confirming that structural information is essential for distinguishing subtle cell states.
\begin{table}[t]
\centering


\small
\begin{tabular}{lccc}
\toprule
\textbf{Model Variants} & \textbf{ACC} & \textbf{NMI} & \textbf{ARI} \\
\midrule
\textbf{\method{}} & \textbf{96.02} & \textbf{89.15} & \textbf{93.10} \\
GE + Pos & 92.12 & 86.33 & 88.55 \\
GE + SP & 80.09 & 71.08 & 67.76 \\
Pos + SP & 27.52 & 12.13 & 5.06 \\
\midrule
Gene Expression (GE) & 79.24 & 70.69 & 67.61 \\
Position (Pos) & 16.68 & 0.73 & 0.23 \\
Shortest Path (SP) & 26.53 & 12.14 & 6.80 \\
\bottomrule
\end{tabular}
\caption{Ablation experiment results of embedding components.}
\label{tab:ablation_embed}
\end{table}

\subsection{Sensitivity Analysis}
To evaluate the robustness of \method{} against hyperparameter variations, we conduct a sensitivity analysis on the Muraro Human Pancreas dataset. We examine neighborhood size $K$ (5--60) and Transformer layers $L$ (1--8). As illustrated in Figure~\ref{fig:sensitivity}, the model demonstrates high stability with respect to $K$, showing only marginal fluctuations across all metrics, indicating that our \method{} is robust to variations in the underlying graph construction. Conversely, performance depends more significantly on network depth $L$. Shallow networks where $L=1$ yield suboptimal scores due to insufficient feature propagation and limited semantic capture. Increasing depth improves results until peaking at $L=6$, where the model effectively balances local and global information. Further extending depth to 8 leads to a slight performance drop, likely due to over-smoothing. Further sensitivity analyses regarding other hyperparameters are detailed in Appendix~\ref{sec:hyperparameter}.